
\subsection{定理 F、G、M、N 与推论 O、P 的证明收束}
\label{subsec:5-4-6}

下面汇集前面各节的结果，
完成定理 F、G、M、N 的证明。

\begin{proof}[定理 M 的证明收束]
考虑由定义 3.8.38 与引理 3.8.39
构造的滤过 $A_\infty$ 代数
$(C(L;\Lambda_{0,\mathrm{nov}}),
\mathfrak m^{\mathbf b})$。
注意
\[
\mathfrak m^{\mathbf b}_1
\equiv\overline{\mathfrak m}_1
\pmod{\Lambda^+_{0,\mathrm{nov}}}.
\]
因此，作为 $\Lambda_{0,\mathrm{nov}}$-模，
$C(L;\Lambda_{0,\mathrm{nov}})$
的典范模型是
$H(L;\Lambda_{0,\mathrm{nov}})$；
记所得结构为
$(H(L;\Lambda_{0,\mathrm{nov}}),
\mathfrak m_{\mathrm{can}}^{\mathbf b})$。
对它应用定义 5.4.11，得到
\[
P_{m+1}^{\mathrm{bos}},
\ldots,P_{m+\ell-1}^{\mathrm{bos}},
P_{m+\ell}^{\mathrm{bos}}
\in\Lambda_{0,\mathrm{nov}}[[s_1,\ldots,s_m]].
\]
这里 $P_{m+\ell}^{\mathrm{bos}}$
是仓西映射沿幺元方向的分量，
即取 $w_\ell=e$。
令
$\mathfrak P_i=P_{m+i}^{\mathrm{bos}}$，
$1\le i\le\ell-1$，
便得到第 1.4 节的
$(\mathfrak P_1,\ldots,\mathfrak P_{\ell-1})$。
于是
\[
\mathcal M_{\mathrm{weak}}^{\mathbf b}(L)
=\bigl(
(\mathfrak P_{L,\mathrm{weak}}^{\mathbf b})^{-1}
(0)(\Lambda^+_{0,\mathrm{nov}})
\bigr)/{\sim}
\]
由命题 5.4.15 的含幺版本推出。
此外，
\[
R(\mathfrak P_{L,\mathrm{weak}}^{\mathbf b};
\Lambda_{0,\mathrm{nov}})
=
\frac{\Lambda_{0,\mathrm{nov}}[[s_1,\ldots,s_m]]}
{(\mathfrak P_1,\ldots,\mathfrak P_{\ell-1})}
\]
正是定理 5.4.12 中的环
$\mathfrak M_{\mathrm{weak}}
(H(L;\Lambda_{0,\mathrm{nov}}),
\mathfrak m_{\mathrm{can}}^{\mathbf b})^{\mathrm{bos}}$。
由定理 4.6.52、命题 5.4.5
和定理 5.4.12，
这个环在连续
$\Lambda_{0,\mathrm{nov}}$-代数同构意义下
与构造中的各种选择无关。
非含幺、非弱的版本同理。
定理 M 证毕。
\end{proof}

\begin{proof}[推论 5.4.21 及定理 F、G 的证明收束]
引理 3.8.74 已在代数对
\[
\bigl(
(C(L^{(1)};\Lambda_{0,\mathrm{nov}}),
\mathfrak m^{\mathbf b}),
(C(L^{(0)};\Lambda_{0,\mathrm{nov}}),
\mathfrak m^{\mathbf b})
\bigr)
\]
上构造同伦含幺的滤过 $A_\infty$ 双模
$(C(L^{(1)},L^{(0)};\Lambda_{0,\mathrm{nov}}),
\mathfrak n^{\mathbf b})$。
由命题 5.2.43，存在含幺滤过
$A_\infty$ 双模 $(D,\mathfrak n)$，
其系数代数对是
\[
\bigl(
(H(L^{(1)};\Lambda_{0,\mathrm{nov}}),
\mathfrak m_{\mathrm{can}}^{\mathbf b}),
(H(L^{(0)};\Lambda_{0,\mathrm{nov}}),
\mathfrak m_{\mathrm{can}}^{\mathbf b})
\bigr),
\]
并有同伦含幺的同伦等价
\begin{equation}
(D,\mathfrak n)
\longrightarrow
(C(L^{(1)},L^{(0)};\Lambda_{0,\mathrm{nov}}),
\mathfrak n^{\mathbf b}),
\tag{5.4.65}\label{eq:5-4-65}
\end{equation}
它覆盖两侧的同伦等价
\[
(H(L^{(i)};\Lambda_{0,\mathrm{nov}}),
\mathfrak m_{\mathrm{can}}^{\mathbf b})
\longrightarrow
(C(L^{(i)};\Lambda_{0,\mathrm{nov}}),
\mathfrak m^{\mathbf b}).
\]
令
$(C_{\mathrm{can}}(L^{(1)},L^{(0)};
\Lambda_{0,\mathrm{nov}}),
\mathfrak n_{\mathrm{can}}^{\mathbf b})$
为 $(D,\mathfrak n)$ 的含幺典范模型
（定理 5.4.18）。
按定义，其底模同构于
$H(D,\overline{\mathfrak n}_{0,0})$。
由 (5.4.65) 可知
\[
C_{\mathrm{can}}(L^{(1)},L^{(0)};
\Lambda_{0,\mathrm{nov}})
\cong
H\bigl(
C(L^{(1)},L^{(0)};\Lambda_{0,\mathrm{nov}}),
\overline{\mathfrak n}_{0,0}^{\mathbf b}
\bigr).
\]
在博特--莫尔斯情形，
$\overline{\mathfrak n}_{0,0}^{\mathbf b}
=\overline{\mathfrak n}_{0,0}
=\pm\partial$，
其中 $\partial$ 是
$\bigcup_hR_h=L^{(1)}\cap L^{(0)}$
上的通常（上）边界算子；
横截相交时该算子为零。
所以
\begin{equation}
C_{\mathrm{can}}(L^{(1)},L^{(0)};
\Lambda_{0,\mathrm{nov}})
=
\bigoplus_{[h,w]\in
\widetilde{\pi}_0(L^{(1)}\cap L^{(0)})}
H^*(R_h;\Theta^-_{R_h})
\otimes_{\mathbb Q}\Lambda_{0,\mathrm{nov}}.
\tag{5.4.66}\label{eq:5-4-66}
\end{equation}
由定理 5.3.1 与命题 5.2.43、5.4.20，
滤过 $A_\infty$ 双模
$(C_{\mathrm{can}}(L^{(1)},L^{(0)};
\Lambda_{0,\mathrm{nov}}),
\mathfrak n_{\mathrm{can}})$
在相应意义下不依赖于各种选择。
这证明了推论 5.4.21，
从而证明定理 F。

接着证明定理 G 的 \textup{(G.2)}。
有
\[
\begin{aligned}
&HF((L^{(1)},b_1),(L^{(0)},b_0);
\Lambda_{\mathrm{nov}})\\
&\quad\cong
H\bigl(
C_{\mathrm{can}}(L^{(1)},L^{(0)};
\Lambda_{0,\mathrm{nov}}),
\delta_{b_1,b_0}
\bigr)
\otimes_{\Lambda_{0,\mathrm{nov}}}
\Lambda_{\mathrm{nov}}.
\end{aligned}
\]
因此 (5.4.66) 给出\translatornote{原书实体页 359
在本秩上界中写
$\Theta_{R_h}$；
同页 (5.4.66) 及第 3.7.5 小节
均使用 $\Theta^-_{R_h}$，
这里保持同一局部系统。
见疑误 A-0341。}
\[
\operatorname{rank}_{\Lambda_{\mathrm{nov}}}
HF((L^{(1)},b_1),(L^{(0)},b_0);
\Lambda_{\mathrm{nov}})
\le
\sum_{h,k}
\operatorname{rank}_{\mathbb Q}
H^k(R_h;\Theta^-_{R_h}),
\]
即所需结论。
\end{proof}

到这里，定理 G 只对
$L^{(1)}$ 与 $L^{(0)}$
干净相交的情形证明完毕；
一般情形将在第 6.5.4 小节处理。

下面收束定理 N 的证明。
考虑 $\mathbf b$ 不必为零的情形，令
\[
\bigl(
\mathfrak D(L^{(1)},L^{(0)};
\Lambda_{0,\mathrm{nov}}),
\mathfrak d
\bigr)
=
\bigl(
\mathfrak M(
C_{\mathrm{can}}(L^{(1)},L^{(0)};
\Lambda_{0,\mathrm{nov}}),
\mathfrak n_{\mathrm{can}}^{\mathbf b}),
\mathfrak Q
\bigr).
\]
右端由命题 5.4.24 定义。
\textup{(N.1)} 随即由引理 5.4.25 得到。
为证明 \textup{(N.2)}，
还需下面的引理；
引理中的各滤过 $A_\infty$ 代数及双模
均假设为有隙的。

\begin{lemma}{5.4.67}\label{lem:5-4-67}
设 $D$、$D'$ 分别为代数对
$(C_1,C_0)$、$(C'_1,C'_0)$
上的滤过 $A_\infty$ 双模，并令
\[
\widetilde D
=D\otimes_{\Lambda_{0,\mathrm{nov}}}
\Lambda_{\mathrm{nov}},
\qquad
\widetilde D'
=D'\otimes_{\Lambda_{0,\mathrm{nov}}}
\Lambda_{\mathrm{nov}}.
\]
设
$\varphi\colon\widetilde D\to\widetilde D'$
是覆盖 $f_i\colon C_i\to C'_i$
的弱滤过 $A_\infty$ 双模同态，
且为同伦等价。
设
$C_{i,\mathrm{can}}\cong C'_{i,\mathrm{can}}$
分别为 $C_i,C'_i$ 的典范模型，
$D_{\mathrm{can}},D'_{\mathrm{can}}$
是相应代数对上的典范滤过
$A_\infty$ 双模，分别与 $D,D'$
同伦等价。
令\translatornote{原书实体页 359
此处把已定义的
$\widetilde D,\widetilde D'$
再次写作与 $\Lambda_{\mathrm{nov}}$
的自身张量积，
而结论及式 (5.4.68)
由映射型别判断均指典范双模的标量扩张；
此处补明下标 $\mathrm{can}$。
见疑误 A-0342。}
\[
\widetilde D_{\mathrm{can}}
=D_{\mathrm{can}}
\otimes_{\Lambda_{0,\mathrm{nov}}}
\Lambda_{\mathrm{nov}},
\qquad
\widetilde D'_{\mathrm{can}}
=D'_{\mathrm{can}}
\otimes_{\Lambda_{0,\mathrm{nov}}}
\Lambda_{\mathrm{nov}}.
\]
则 $\widetilde D_{\mathrm{can}}$
与 $\widetilde D'_{\mathrm{can}}$
在恒等底同态上同伦等价。
\end{lemma}
\begin{proof}
取同伦等价
$g_i\colon C_{i,\mathrm{can}}\to C_i$、
$g'_i\colon C'_{i,\mathrm{can}}\to C'_i$，
以及它们的同伦逆
$g_i^{-1},g_i'^{-1}$。
取
$\phi\colon D_{\mathrm{can}}\to D$、
$\phi'\colon D'_{\mathrm{can}}\to D'$
分别覆盖 $g_i,g'_i$
的同伦等价，
及覆盖相应同伦逆的
$\phi^{-1},\phi'^{-1}$。
\translatornote{原书实体页 360
把 $\phi,\phi'$ 均称为覆盖
$f_i$，其逆则覆盖 $f_i^{-1}$；
但各映射两端分别同属无撇或带撇系数代数，
只有覆盖 $g_i,g'_i$
才能使后续复合的底同态有型别。
见疑误 A-0343。}
于是
\begin{equation}
\bigl(
\phi'^{-1}\otimes_{\Lambda_{0,\mathrm{nov}}}
\Lambda_{\mathrm{nov}}
\bigr)
\circ\varphi\circ
\bigl(
\phi\otimes_{\Lambda_{0,\mathrm{nov}}}
\Lambda_{\mathrm{nov}}
\bigr)
:
\widetilde D_{\mathrm{can}}
\longrightarrow
\widetilde D'_{\mathrm{can}}
\tag{5.4.68}\label{eq:5-4-68}
\end{equation}
是覆盖
$g_i'^{-1}\circ f_i\circ g_i$
的同伦等价。
可以选取 $g_i,g'_i$，
使这个底同态同伦于恒等同态。
由命题 5.2.43，
可将 (5.4.68) 形变为
覆盖恒等同态的同伦等价。
\end{proof}

由引理 5.4.67 和定理 5.3.14
可得定理 N 的 \textup{(N.2)}；
定理 N 证毕。

\begin{remark}{5.4.69}\label{rem:5-4-69}
在引理 5.4.67 的情形，
$\widetilde D_{\mathrm{can}}$
与 $\widetilde D'_{\mathrm{can}}$
一般并不同构。
事实上，在定理 F 的情形，
\[
\operatorname{rank}_{\Lambda_{0,\mathrm{nov}}}
C_{\mathrm{can}}(L^{(1)},L^{(0)};
\Lambda_{0,\mathrm{nov}})
\]
等于 $L^{(1)}\cap L^{(0)}$
的贝蒂数之和；
当用两个哈密顿微分同胚
分别移动 $L^{(0)},L^{(1)}$ 时，
该数并不是不变量。
\end{remark}

现在证明推论 O、P。
设 $(M,\omega)$ 为满足 $c_1(M)=0$
的辛流形，
$L$ 为满足
$\mu_L:H_2(M,L)\to\mathbb Z$
为零的拉格朗日子流形；
对 $L^{(1)},L^{(0)}$
也作同样假设。

构造滤过 $A_\infty$ 代数
$(C(L;\Lambda_{0,\mathrm{nov}}),\mathfrak m)$
时，生成元 $e$ 总以
$e^{\mu_L(\beta)/2}$ 的形式出现。
因此在当前情形，可在
$\Lambda_0$ 上定义这个滤过
$A_\infty$ 代数
（见第 1 章约定 \textup{(Conv.5)}）。
定义 $\Lambda_+$ 时要求各指数
$\lambda_i>0$。
第 3--5 章的代数论证
用 $\Lambda_0$ 代替
$\Lambda_{0,\mathrm{nov}}$
后仍成立。
故若
$\mathbf b\in H^2(M;\Lambda_+)$，
便有典范模型
$(H(L;\Lambda_0),
\mathfrak m_{\mathrm{can}}^{\mathbf b})$。
\translatornote{原书实体页 360
在此把 bulk 类写成
$H^2(L;\Lambda_+)$ 中的类；
但定义 3.8.38 与推论 O
都以 $H^2(M;\Lambda_+)$
中的类作为内部标记点参数。
见疑误 A-0352。}

为在此情形下实行第 5.4.1 小节的构造，
定义仓西映射，令
$m=\operatorname{rank}H^1(L;\mathbb Q)$。
注意，第 5.4.1 小节的 $m$
则是 $\operatorname{rank}H^{\mathrm{odd}}
(L;\mathbb Q)$。
取 $H^1(L;\Lambda_0)
=\Lambda_0^{\oplus m}$
的一组基
$v_1,\ldots,v_m$。
类似地，令
$\ell=\operatorname{rank}H^2(L;\mathbb Q)$，
取 $H^2(L;\Lambda_0)
=\Lambda_0^{\oplus\ell}$
的一组基
$w_1,\ldots,w_\ell$，
并以 $\pi_i$
表示投影到第 $i$ 个分量。
仿照 (5.4.9) 定义\translatornote{原书实体页 360
在 (5.4.70) 的值域仍写
$\Lambda_{0,\mathrm{nov}}[[s_1,\ldots,s_m]]$；
本段已改用 $\Lambda_0$，
且推论 P 的坐标环正是
$\Lambda_0[[s_1,\ldots,s_m]]$。
见疑误 A-0344。}
\begin{equation}
\begin{aligned}
P_i(s_1,\ldots,s_m)
&=(\pi_i\circ\mathfrak m^{\mathbf b})
\bigl(\exp(s_1v_1+\cdots+s_mv_m)\bigr)\\
&\in\Lambda_0[[s_1,\ldots,s_m]].
\end{aligned}
\tag{5.4.70}\label{eq:5-4-70}
\end{equation}
于是得到
\[
\mathfrak P_L^{\mathbf b}
=(P_1,\ldots,P_\ell):
H^1(L;\Lambda_+)
\longrightarrow H^2(L;\Lambda_+).
\]
利用下述引理，
仿照定理 M 的证明给出\translatornote{原书实体页 360
写作 “Theorem P”，
但不存在该定理；
此处的模空间与仓西映射零点之等式
正对应紧前定理 M 的论证。
见疑误 A-0350。}
\[
\mathcal M^{\mathbf b}(L;\Lambda_+)
=
(\mathfrak P_L^{\mathbf b})^{-1}
(0)(\Lambda_+).
\]

\begin{lemma}{5.4.71}\label{lem:5-4-71}
设 $C$ 是 $\Lambda_0$ 上的
含幺滤过 $A_\infty$ 代数，
且 $C^0=\Lambda_0e$，
其中 $e$ 是幺元。
那么，在
\[
\widehat{\mathcal M}(C)
=\{b\in C^1\mid
\mathfrak m(e^b)=0\}
\]
上的规范等价关系 $\sim$
是平凡的。
\end{lemma}
\begin{proof}
设 $b_1,b_0\in\widehat{\mathcal M}(C)$，
且 $b_1\sim b_0$。
由命题 4.3.9，
存在
$(b_1,c,b_0)\in C^{[0,1]}$，
使
\[
b_1-b_0
=\mathfrak m(e^{b_1}ce^{b_0}).
\]
因为 $\deg b_i=1$，
所以 $\deg c=0$；
由假设可写 $c=\rho e$，
其中 $\rho\in\Lambda_0$。
含幺性使
$\mathfrak m(e^{b_1}ce^{b_0})=0$，
于是 $b_1=b_0$。
\end{proof}

接着证明 \textup{(O.4)}。
设正整数 $m$ 满足
$m[\omega]\in H^2(M,L;\mathbb Z)$。
在构造
$(C(L;\Lambda_{0,\mathrm{nov}}),\mathfrak m)$
时，参数 $T$ 只以
$T^{\omega\cap\beta}$ 的形式出现，
其中 $\beta\in\pi_2(M,L)$。
因此系数环可缩减为
$\mathbb Q[[T^{1/m}]]$。

现考虑两个马斯洛夫指标均为零的
拉格朗日子流形
$L^{(1)},L^{(0)}$，
以证明 \textup{(O.1)}、
\textup{(O.2)}。
这正是第 5.1.3 小节
注 5.1.18 所讨论的要点。
下面沿用第 2.2.1 小节的记号
证明一个引理。

\begin{lemma}{5.4.72}\label{lem:5-4-72}
若
$\mu_{L^{(i)}}:
H_2(M,L^{(i)})\to\mathbb Z$
为零，则
\[
I_\mu:
\pi_1\bigl(
\Omega(L^{(0)},L^{(1)};\ell_0)
\bigr)\longrightarrow\mathbb Z
\]
也为零。
\end{lemma}
\begin{proof}
令 $\operatorname{Lag}(M)\to M$
为拉格朗日格拉斯曼流形丛，
其在 $p\in M$ 处的纤维是
$\Lambda(T_pM)$
（见第 2.1.1 小节）。
由 \textbf{[Sei00, Fuk02II]}
及 $c_1(M)=0$，
存在一个 $\mathbb Z$-覆盖
\[
\widetilde{\operatorname{Lag}}(M)
\longrightarrow\operatorname{Lag}(M),
\]
使
$\widetilde{\operatorname{Lag}}(M)\to M$
的每个纤维是
$\Lambda(T_pM)$ 的万有覆盖。

限制在 $L^{(i)}$ 上时，
$\operatorname{Lag}(M)$
有截面
$s_{L^{(i)}}(p)
=T_pL^{(i)}$。
把它提升为
$\widetilde s_{L^{(i)}}:
L^{(i)}\to
\widetilde{\operatorname{Lag}}(M)$
的障碍，
是 $\mu_{L^{(i)}}\in
H^2(M,L^{(i)})$
在 $H^1(L^{(i)})$ 中的像
（见 \textbf{[Fuk02II]}）。
按假设此像为零，
故这样的提升存在。

取映射
$C:S^1\times[0,1]\to M$，
满足
$C(S^1\times\{i\})\subset L^{(i)}$。
它代表
$\pi_1(\Omega(L^{(0)},L^{(1)};\ell_0))$
中的一个元素。
原书以“纤维连通”为由断言，
沿 $C$ 拉回的
$\widetilde{\operatorname{Lag}}(M)$
是平凡丛；
\translatornote{原书实体页 361
括注只说“因为纤维连通”；
这单独不足以推出圆柱上的纤维丛平凡。
此处保留原书的断言与后续推理，
但不把该括注当作完整证明。
见疑误 A-0351。}
故 $\widetilde s_{L^{(i)}}$
给出从
$S^1\times\{i\}$
到 $\widetilde{\operatorname{Lag}}(M)$
的映射。
由于
\[
\pi_1(\widetilde{\operatorname{Lag}}(M))
\longrightarrow\pi_1(M)
\]
是同构，
$\widetilde s_{L^{(0)}}$
与 $\widetilde s_{L^{(1)}}$
同伦。
于是存在沿 $C$ 的拉回丛的截面
$\widetilde s$，
在两条边界上分别限制为
$\widetilde s_{L^{(i)}}$；
投影后得到
$\operatorname{Lag}(M)$
的拉回丛截面 $s$，
在边界上等于
$s_{L^{(i)}}$。
依 $I_\mu$ 的定义，
这说明 $I_\mu(C)=0$。
\end{proof}

再假设 $L^{(1)}$
与 $L^{(0)}$ 干净相交，
并记
$L^{(0)}\cap L^{(1)}
=\bigcup_hR_h$。
对道路空间
$\Omega(L^{(0)},L^{(1)})$
的每个连通分支，
选取一条通向某个
$R_{h_0}$ 的常值基道路 $\ell_0$。
引理 5.4.72 表明，
$\mu([h,w];\ell_0)$
与 $w$ 无关，
仅依赖于 $h$。
令\translatornote{原书实体页 361
此式漏去按交集分量 $h$ 的直和，
并把局部系统写成
$\Theta_{R_h}$；
第 3.7.5 小节的同一复形
有分量直和及
$\Theta^-_{R_h}$。
这里恢复两者并取能量完备化。
见疑误 A-0346、A-0348。}
\[
\widehat{CF}_{BM}
(L^{(1)},L^{(0)};\Lambda_0)
=
\widehat{\bigoplus}_{h}
\left(
C(R_h;\Theta^-_{R_h})
\otimes_{\mathbb Q}\Lambda_h
\otimes_{\mathbb Q}\Lambda_0
\right).
\]
这与第 3.7.5 小节
$\Lambda_{0,\mathrm{nov}}$ 系数的
定义相似。
在此空间上修改 (3.7.19)
以定义等价关系：
\[
S\otimes[h,w]\otimes T^\lambda
\sim
S'\otimes[h',w']\otimes T^{\lambda'}
\]
当且仅当
\begin{align}
h&=h',
\tag{5.4.73.1}\label{eq:5-4-73-1}\\
S&=S',
\tag{5.4.73.2}\label{eq:5-4-73-2}\\
\lambda+\int_w\omega
&=\lambda'+\int_{w'}\omega.
\tag{5.4.73.3}\label{eq:5-4-73-3}
\end{align}
这里不再有与 (3.7.19.3)
相对应的条件。
将商空间
$\widehat{CF}_{BM}
(L^{(1)},L^{(0)};\Lambda_0)/{\sim}$
的非负能量部分记作
$C(L^{(1)},L^{(0)};\Lambda_0)$
（见第 3.7.4、3.7.5 小节）。
由 $\mu([h,w];\ell_0)$
与 $w$ 无关，可得\translatornote{原书实体页 362
此式把局部系数写成
$\Theta_{R_h}$，
并把张量底环写成 $\mathbb R$；
上一页的定义与第 3.7.5 小节
均使用 $\Theta^-_{R_h}$
及 $\mathbb Q$。
见疑误 A-0346、A-0349。}
\[
C(L^{(1)},L^{(0)};\Lambda_0)
\cong
\bigoplus_h
C(R_h;\Theta^-_{R_h})
\widehat\otimes_{\mathbb Q}\Lambda_0.
\]
此外，可在其上定义次数：
\[
\deg(S\otimes[h,w]\otimes T^\lambda)
=\deg S+\mu([h,w];\ell_0).
\]
现在只需从公式中去掉 $e$，
便可在 $\Lambda_0$ 上
重复第 3.7 节滤过
$A_\infty$ 双模的整个构造。
与定理 G 的证明相同，
这给出 \textup{(O.1)}、
\textup{(O.2)}、
\textup{(O.3)}。
推论 O 证毕。
\translatornote{原书实体页 362
末句误作 “Corollary R”；
本段证明的是推论 O，
而 R 在第 1.4 节是猜想。
见疑误 A-0347。}

\begin{remark}{5.4.74}\label{rem:5-4-74}
若 $L^{(i)}$ 在
\textbf{[Fuk03I]} 的意义下有理
（见注 6.2.4），
则可把
$C(L^{(1)},L^{(0)};\Lambda_0)$
以及弗洛尔上同调的系数环
缩减为
$\mathbb Q[[T^{1/m}]][T^{-1}]$。
有理性并不蕴含
对每个 $B\in\pi(p,q)$
都有 $\omega(B)\in\mathbb Q$
（记号见定义 2.3.7）。
所以还需修改边界算子，
使它在
$\mathbb Q[[T^{1/m}]][T^{-1}]$
上有定义；
修改方法见 \textbf{[Fuk03I]}。
\end{remark}
